\section*{Bab \ref{ch:g9:statproba} --- Statistika dan Peluang}

\begin{solution}{exo:g9:statproba:1}
\omterm{def:g9:statproba:indicators}{Rata-rata}: $\dfrac{14 + 9 + 17 + 9 + 12 + 11 + 12}{7} = \dfrac{84}{7} = 12$.

\omterm{def:g9:statproba:indicators}{Median}: terurut menjadi $9, 9, 11, 12, 12, 14, 17$; nilai tengahnya
(yang ke-$4$) adalah $12$.

\omterm{def:g9:statproba:indicators}{Jangkauan}: $17 - 9 = 8$.
\end{solution}

\begin{solution}{exo:g9:statproba:2}
\omterm{def:g9:statproba:indicators}{Rata-rata}: $\dfrac{19 + 21 + 24 + 18 + 22 + 25 + 19}{7} = \dfrac{148}{7}
\approx 21.1$ derajat.

\omterm{def:g9:statproba:indicators}{Median}: terurut menjadi $18, 19, 19, 21, 22, 24, 25$; nilai ke-$4$
adalah $21$ derajat.
\end{solution}

\begin{solution}{exo:g9:statproba:3}
\omterm{def:g9:statproba:indicators}{Rata-rata}:
\[
\frac{2 \times 0 + 4 \times 1 + 3 \times 2 + 1 \times 3}{10}
= \frac{0 + 4 + 6 + 3}{10} = \frac{13}{10} = 1.3
\text{ gol tiap pertandingan.}
\]
\omterm{def:g9:statproba:indicators}{Median}: deret terurutnya adalah $0,0,1,1,1,1,2,2,2,3$; nilai ke-$5$ dan
ke-$6$ keduanya $1$: jadi \omterm{def:g9:statproba:indicators}{mediannya} $1$.
\end{solution}

\begin{solution}{exo:g9:statproba:4}
$\P(3) = \frac16$. \quad
$\P(\text{genap}) = \frac36 = \frac12$. \quad
$\P(\text{sedikitnya } 3) = \P(\{3,4,5,6\}) = \frac46 = \frac23$. \quad
$\P(\text{bukan } 6) = 1 - \frac16 = \frac56$.
\end{solution}

\begin{solution}{exo:g9:statproba:5}
$\P(\text{putih}) = \frac{5}{12}$.

$\P(\text{merah atau hijau}) = \frac{4 + 3}{12} = \frac{7}{12}$.

$\P(\text{bukan putih}) = 1 - \frac{5}{12} = \frac{7}{12}$.

Kedua yang terakhir sama: ``bukan putih'' \emph{adalah} kejadian
``merah atau hijau''.
\end{solution}

\begin{solution}{exo:g9:statproba:6}
Semua $12$ hasilnya (sisi koinnya, nilai dadunya) sama-sama
mungkin.

Sisi gambar dan angka $6$: satu hasil dari $12$: jadi $\frac{1}{12}$.

Sisi gambar atau angka $6$: hasil dengan sisi gambar ($6$ buah)
ditambah hasil (sisi angka, $6$): jadi $7$ hasil, yaitu
$\frac{7}{12}$. Membilang tiap kejadiannya lalu mengurangkan
tumpang tindihnya memberi hasil yang sama:
$\frac{6}{12} + \frac{2}{12} - \frac{1}{12} = \frac{7}{12}$.
\end{solution}

\begin{solution}{exo:g9:statproba:7}
Pengambilan pertama: merah $\frac25$, hitam $\frac35$. Pengambilan
kedua tanpa pengembalian: sesudah merah, tersisa $1$ merah dan $3$
hitam di antara $4$; sesudah hitam, tersisa $2$ merah dan $2$ hitam di
antara $4$.

Dua hitam: $\frac35 \times \frac24 = \frac{6}{20} = \frac{3}{10}$.

Warna yang berbeda (MH atau HM):
$\frac25 \times \frac34 + \frac35 \times \frac24
= \frac{6}{20} + \frac{6}{20} = \frac{12}{20} = \frac35$.
\end{solution}

\begin{solution}{exo:g9:statproba:8}
\emph{1.} $9 \times 14 = 126$ angka.

\emph{2.} \omterm{def:g9:statproba:indicators}{Rata-rata} $15$ atas $10$ pertandingan berarti $150$ angka
seluruhnya: jadi ia harus mencetak $150 - 126 = 24$ angka.
\end{solution}

\begin{solution}{exo:g9:statproba:9}
\emph{1.} Di antara $36$ pasangan yang sama-sama mungkin, dobelnya
adalah $(1,1), (2,2), \dots, (6,6)$: enam buah, jadi
$\P(\text{dobel}) = \frac{6}{36} = \frac16$.

\emph{2.} Dua permainan yang saling bebas: pohon dengan menang
($\frac16$) atau kalah ($\frac56$) pada tiap tingkatnya. \omterm{def:g9:statproba:probability}{Peluang} tidak
pernah menang: $\frac56 \times \frac56 = \frac{25}{36}$. Jadi
\[
\P(\text{sedikitnya satu kemenangan}) = 1 - \frac{25}{36} = \frac{11}{36}
\approx 0.31 .
\]
\end{solution}

\begin{solution}{pb:g9:statproba:1}
\textbf{1.} Tiap lemparannya punya $5$ hasil bukan enam dari $6$;
tabel $6 \times 6$ menunjukkan $5 \times 5 = 25$ hasil bebas angka
enam dari $36$: jadi $P(\text{tanpa enam}) = \frac{25}{36}$, sehingga
$P(\text{sedikitnya satu}) = 1 - \frac{25}{36} = \frac{11}{36}
\approx 0.31$.

\textbf{2.} $\frac13 = \frac{12}{36} \neq \frac{11}{36}$.
Penjumlahannya membilang hasil ``enam lalu enam'' sekali pada tiap
$\frac16$: yaitu sekali terlalu sering. \omterm{def:g9:statproba:probability}{Peluang} kejadian yang dapat
terjadi \emph{bersama} tidak begitu saja dijumlahkan.

\textbf{3.} \omterm{def:g1:addition:def}{Jumlah} $7$: keenam hasilnya $(1,6), (2,5), (3,4),
(4,3), (5,2), (6,1)$: jadi $P = \frac{6}{36} = \frac16$. \omterm{def:g1:addition:def}{Jumlah} $2$:
hanya $(1,1)$: jadi $\frac{1}{36}$. Angka tujuh punya cara terbanyak
untuk terjadi --- yaitu bagian tengah yang gemuk pada tabel \omterm{def:g1:addition:def}{jumlahnya}.

\textbf{4.} \omterm{def:g1:addition:def}{Jumlah} $9$: $(3,6), (4,5), (5,4), (6,3)$ --- yaitu empat
hasil, $\frac{4}{36}$. \omterm{def:g1:addition:def}{Jumlah} $10$: $(4,6), (5,5), (6,4)$ ---
tiga hasil, $\frac{3}{36}$. Angka sembilan lebih mungkin: bilangannya
yang memutuskan, bukan banyaknya ``cara menulis'' \omterm{def:g1:addition:def}{jumlahnya} dengan
sisi dadu yang tidak berurutan.

\textbf{5.} Tiap tingkat pohonnya mengalikan \omterm{def:g9:statproba:probability}{peluang} tanpa angka enam
dengan $\frac56$, apa pun yang terjadi sebelumnya: jadi sesudah $n$
lemparan, $P(\text{tanpa enam sama sekali}) = \left(\frac56\right)^n$,
dan aturan komplemennya memberi rumusnya.

\textbf{6.} $\left(\frac56\right)^4 = \frac{625}{1296} \approx
0.482$, jadi Chevalier menang dengan \omterm{def:g9:statproba:probability}{peluang}
$1 - 0.482 = 0.518$ --- yaitu keunggulan yang tenang dan mantap
sebesar hampir $2\,\%$ tiap permainan: taruhan yang sangat baik, dan
taruhan itu membayar keretanya.

\textbf{7.} Pada $7$ lemparan aturannya memberi $\frac76$ --- yaitu
\omterm{def:g9:statproba:probability}{peluang} yang lebih besar daripada $1$, jadi omong kosong. Itulah
pembilangan ganda pada pertanyaan 2, yang menumpuk: keberhasilan
lemparannya tumpang tindih, dan menjumlahkan \omterm{def:g9:statproba:probability}{peluangnya} membilang
tumpang tindihnya berulang kali.

\textbf{8.} $1 - \left(\frac{35}{36}\right)^{24} \approx
1 - 0.5086 = 0.491$: jadi di bawah setengah. Dengan bertaruh uang
seimbang pada kejadian $49.1\,\%$, Chevalier kalah kira-kira $2$
permainan tiap $100$ secara \omterm{def:g9:statproba:indicators}{rata-rata} --- perlahan, secara misterius
(baginya), dan tidak terhindarkan.

\textbf{9.} $1 - \left(\frac{35}{36}\right)^{25} \approx
1 - 0.494 = 0.506$: jadi dengan $25$ lemparan taruhannya berbalik
menguntungkan. Satu lemparan memisahkan Chevalier dari keuntungan.

\textbf{10.} Mengalikan sebuah \omterm{def:g9:statproba:probability}{peluang} dengan banyaknya percobaan
membilang keberhasilan yang tumpang tindih beberapa kali, dan
akhirnya menghasilkan jawaban mustahil di atas $1$. Jalan yang benar
selalu lewat komplemennya: rangkailah \omterm{def:g9:statproba:probability}{peluang} kegagalannya
dengan perkalian, lalu kurangkan dari $1$.

\textbf{11.} \omterm{def:g9:statproba:indicators}{Rata-rata}: $\frac{9 \times 2\,000 + 20\,000}{10} =
\frac{38\,000}{10} = 3\,800$ euro. \omterm{def:g9:statproba:indicators}{Median}: gaji tengahnya
keduanya $2\,000$: jadi \omterm{def:g9:statproba:indicators}{mediannya} $2\,000$ euro. Iklannya
sebaiknya mengutip \omterm{def:g9:statproba:indicators}{mediannya} --- sembilan dari sepuluh pekerjanya
tidak pernah melihat apa pun yang mendekati $3\,800$; \omterm{def:g9:statproba:indicators}{rata-ratanya}
terseret naik oleh satu pencilan
(\cref{ex:g9:statproba:weighted} sudah mengingatkan soal bobot).

\textbf{12.} Misalnya $2, 8, 15, 17, 18$: \omterm{def:g9:statproba:indicators}{mediannya} $15$
(yaitu nilai tengahnya), \omterm{def:g9:statproba:indicators}{rata-ratanya} $\frac{60}{5} = 12$. \omterm{def:g9:statproba:indicators}{Rata-rata}
di bawah \omterm{def:g9:statproba:indicators}{mediannya} mengkhianati ekor nilai yang \emph{rendah} yang
menarik \omterm{def:g9:statproba:indicators}{rata-ratanya} turun, sedangkan separuh atasnya duduk
tinggi.

\textbf{13.} Ketiga faktor pertamanya:
$1 \times \frac{364}{365} \times \frac{363}{365} \approx
0.992$. Tiap murid baru harus menghindari satu tanggal terpakai lagi,
jadi faktornya mengerut: $\frac{343}{365} \approx 0.94$ untuk murid
ke-$23$. \omterm{def:g2:mult:def}{Hasil kali} seluruhnya jatuh ke kira-kira $0.493$ --- jadi
\[
P(\text{sedikitnya satu ulang tahun yang sama}) \approx 0.507 :
\]
jadi lebih baik daripada seimbang, pada kelas berisi hanya $23$ murid.
Gerak hati berkata ``$23$ orang, $365$ hari, tidak mungkin'';
matematikanya berkata ``lemparkan koin''.

\textbf{14.} Ada $\frac{23 \times 22}{2} = 253$ pasang. Tiap pasangnya
adalah kesempatan baru untuk sebuah kebetulan, dan $253$ kesempatan
pada $\frac{1}{365}$ masing-masing bukan lagi perkara kecil: jadi
kejutannya larut begitu kita membilang pasangan, bukan orang.

\textbf{15.} Contoh tata caranya: lemparkan dadu dalam blok berisi
empat lemparan, sambil mencatat apakah tiap bloknya memunculkan angka
enam, untuk $100$ blok; atau kumpulkan daftar ulang tahun banyak kelas
yang berisi kira-kira $23$ murid lalu catat berapa bagian yang punya
kebetulan. \omterm{def:g7:stats:frequency}{Frekuensi relatif} yang teramati akan bergoyang, tetapi
ketika banyaknya pengulangan bertambah, \omterm{def:g7:stats:frequency}{frekuensi} itu seharusnya
menetap makin dekat dengan \omterm{def:g9:statproba:probability}{peluang} yang dihitung
($0.518$; $0.507$) --- jadi \omterm{def:g7:stats:frequency}{frekuensi relatif} menuju \omterm{def:g9:statproba:probability}{peluang}: itulah
hukum bilangan besar, yang dinyatakan di sini sebagai harapan dan
dibuktikan di buku Sekolah Menengah.

\end{solution}
